\section{The discrete survival constant \texorpdfstring{\(A(p)\)}{A(p)}}
\label{sec:Ap}

For \(0<p<\tfrac12\), every late conditional moment in
\cref{sec:qsd} depends on
\[
  A(p)=\lim_{n\to\infty}\frac{S_n}{(2p)^n}.
\]
The recursion for \(S_n\) makes numerical evaluation straightforward away from
criticality, but it does not produce an elementary formula.  This section first
extracts what can be proved directly: existence, exact product and series
representations, bounds and a sharp near-critical asymptotic.  Only then do we
turn to the question of closed form and to practical evaluation.

The endpoint \(p=0\) is excluded from the defining ratio, which would be \(0/0\)
after the first generation.  We shall use \(A(0)=\tfrac12\) only for the continuous
extension established below.  At the other endpoint \(A(\tfrac12)=0\) is likewise
an extension from the subcritical interval.

\subsection{Logistic iteration and an infinite product}
\label{sec:product}

Put
\begin{equation}
  m=2p,
  \qquad
  w_n=\frac{S_n}{2}.
  \label{eq:logistic-change}
\end{equation}
Then \eqref{SnIter} becomes the logistic iteration
\begin{equation}
  w_{n+1}=f_m(w_n),
  \qquad
  f_m(w)=mw(1-w),
  \qquad
  w_0=\frac12,
  \label{eq:logistic}
\end{equation}
in the parameter range \(0<m<1\).  The origin is attracting throughout this
range; the period-doubling behaviour of the logistic family at larger \(m\) is not
involved.

Define \(x_n=w_n/m^n\).  From \eqref{eq:logistic},
\begin{equation}
  x_{n+1}=x_n(1-w_n),
  \qquad
  x_n=\frac12\prod_{k=0}^{n-1}(1-w_k).
  \label{eq:xn-product}
\end{equation}
Since \(A(p)=2\lim_nx_n\), this gives the exact product
\begin{equation}
  A(p)=\prod_{k=0}^{\infty}(1-w_k)
  =\frac12\prod_{k=1}^{\infty}
  \left(1-\frac{S_k}{2}\right).
  \label{eq:prodFormula}
\end{equation}
The convergence is not merely formal.  For \(0<w_n<1\),
\(w_{n+1}<mw_n\), so \(w_n\leq m^n/2\) and
\(\sum_nw_n<\infty\).  The infinite product therefore converges to a strictly
positive value.  This proves existence of \(A(p)\) and the asymptotic
\(S_n\sim A(p)m^n\).

Equation \eqref{eq:xn-product} also corrects a possible computational
misunderstanding.  The raw ratios are themselves the partial products:
\begin{equation}
  \frac{S_n}{m^n}=\prod_{k=0}^{n-1}
  \left(1-\frac{S_k}{2}\right).
  \label{eq:ratio-is-product}
\end{equation}
They decrease monotonically to \(A(p)\), and hence every finite ratio is already
a rigorous upper bound.

\subsection{An exact series and two-sided bounds}
\label{sec:Abounds}

The reciprocal constant has an additive representation that exposes the
near-critical scale.  Let \(v_n=m^n/S_n\).  Using
\(S_{n+1}=mS_n(1-S_n/2)\),
\begin{align}
  v_{n+1}-v_n
  &=\frac{m^n}{S_n}
    \left(\frac{1}{1-S_n/2}-1\right)
    =\frac{m^n}{2-S_n}.
  \label{eq:vn-increment}
\end{align}
Since \(v_0=1\) and \(v_n\to1/A(p)\), summation gives the following exact
identity.

\begin{proposition}[Series representation]
For \(0<p<\tfrac12\),
\begin{equation}
  \frac{1}{A(p)}
  =1+\sum_{n=0}^{\infty}\frac{(2p)^n}{2-S_n}.
  \label{eq:Aseries}
\end{equation}
\end{proposition}

The survival probability lies in \((0,1]\), so the denominator in every summand
lies in \([1,2)\).  This simple observation gives bounds that remain informative
at both ends of the subcritical interval.

\begin{proposition}[Two-sided bounds]
Let \(\varepsilon=1-2p\).  For \(0<p<\tfrac12\),
\begin{equation}
  \frac{\varepsilon}{1+\varepsilon}
  <A(p)<
  \frac{2\varepsilon}{1+2\varepsilon},
  \label{eq:Abounds}
\end{equation}
or, in terms of \(p\),
\begin{equation}
  \frac{1-2p}{2(1-p)}
  <A(p)<
  \frac{2(1-2p)}{3-4p}.
  \label{eq:Aboundsexplicit}
\end{equation}
In particular,
\begin{equation}
  \frac12A_{\mathrm c}(p)<A(p)<A_{\mathrm c}(p).
  \label{eq:Ac-comparison-bound}
\end{equation}
\end{proposition}

\begin{proof}
For every \(n\),
\[
  \frac12m^n<\frac{m^n}{2-S_n}\leq m^n,
\]
with strictness somewhere on both sides.  Substitution into
\eqref{eq:Aseries} and summation of the geometric series yield
\[
  1+\frac{1}{2\varepsilon}
  <\frac1{A(p)}<
  1+\frac1{\varepsilon}.
\]
Inverting proves \eqref{eq:Abounds}.  The first expression in
\eqref{eq:Aboundsexplicit} is \(A_{\mathrm c}(p)/2\), while its upper bound is
strictly below \(A_{\mathrm c}(p)=2\varepsilon/(1+\varepsilon)\).
\end{proof}

Binary reproduction supplies a second constraint.  For every \(n\geq1\), a
surviving population has even size and is therefore at least two.  Hence
\begin{equation}
  \E[Z_n\mid Z_n>0]\geq2,
  \qquad
  A(p)\leq\frac12.
  \label{eq:parity-bound}
\end{equation}
The inequality is strict for \(p>0\), but the lower bound in
\eqref{eq:Aboundsexplicit} tends to \(1/2\) as \(p\downarrow0\).  It follows by
squeezing that
\begin{equation}
  A(p)\longrightarrow\frac12
  \qquad (p\downarrow0).
  \label{eq:A-at-zero}
\end{equation}
This limit defines the endpoint value \(A(0)\).

\subsection{The near-critical limit}
\label{sec:asymptotic}

The bounds determine the order of \(A(p)\) near criticality but leave a factor of
two unresolved.  The exact series fixes that factor.  The argument is short, but
the uniform comparison it uses is important: a pointwise approximation of
\(S_n\) would not by itself justify summation over \(n\) of order
\(1/\varepsilon\).

\begin{theorem}[Near-critical survival amplitude]
As \(p\uparrow\tfrac12\),
\begin{equation}
  A(p)\sim2(1-2p).
  \label{eq:Aasym}
\end{equation}
Equivalently, the quasi-stationary mean satisfies
\(
  1/A(p)\sim[2(1-2p)]^{-1}.
\)
\end{theorem}

\begin{proof}
Write \(m=1-\varepsilon\), and make the dependence of the survival sequence on
\(m\) explicit as \(S_n(m)\).  The map
\(mS(1-S/2)\) is increasing in \(m\) and in \(S\in[0,1]\); induction therefore
gives
\begin{equation}
  0\leq S_n(m)\leq S_n(1)
  \qquad (0<m\leq1).
  \label{eq:Sn-critical-domination}
\end{equation}
By \eqref{twoovrn_sn}, \(S_n(1)\to0\).  Given \(\eta>0\), choose \(N\) such
that \(S_N(1)<\eta\).  Monotonicity in \(n\) and
\eqref{eq:Sn-critical-domination} then imply
\[
  \frac12
  \leq\frac{1}{2-S_n(m)}
  \leq\frac{1}{2-\eta}
  \qquad(n\geq N).
\]
Multiply \eqref{eq:Aseries} by \(\varepsilon\).  The initial \(1\) and the first
\(N\) summands contribute \(o(1)\), whereas
\(
  \varepsilon\sum_{n=N}^{\infty}m^n=m^N\to1.
\)
Thus
\[
  \frac12
  \leq\liminf_{\varepsilon\downarrow0}\frac{\varepsilon}{A(p)}
  \leq\limsup_{\varepsilon\downarrow0}\frac{\varepsilon}{A(p)}
  \leq\frac{1}{2-\eta}.
\]
Letting \(\eta\downarrow0\) gives
\(
  \varepsilon/A(p)\to1/2
\), which is \eqref{eq:Aasym}.
\end{proof}

The continuum Riccati equation in \cref{sec:riccati} reaches the same leading
scale by matching its exponential tail to the critical \(2/n\) decay.  It is a
useful interpretation, but the preceding proof shows exactly why the matching is
valid.  Numerically, \(A(p)/[2(1-2p)]\) approaches one slowly.  With
\(\varepsilon=1-2p\), high-precision calculations fit the deficit as
\begin{equation}
  1-\frac{A(p)}{2\varepsilon}
  \approx \varepsilon\bigl[\log(1/\varepsilon)+0.78\bigr].
  \label{eq:A-empirical-remainder}
\end{equation}
This is an empirical rate observation, not a proved remainder estimate.

\subsection{Discrete and continuous time}
\label{sec:compare}

The continuous-time constant from \eqref{eq:Acts} is
\[
  A_{\mathrm c}(p)=\frac{1-2p}{1-p}.
\]
At strong subcriticality, \eqref{eq:A-at-zero} gives
\(A(p)\to1/2\), whereas \(A_{\mathrm c}(p)\to1\).  This factor of two is a
state-space effect.  A late survivor in the discrete binary process must contain
at least two individuals; in the continuous-time process it may be a single
individual whose death clock has not yet fired.  Their limiting conditional means
therefore approach two and one, respectively.

Near criticality, by contrast,
\[
  A_{\mathrm c}(p)
  =\frac{2\varepsilon}{1+\varepsilon}
  \sim2\varepsilon
  \sim A(p).
\]
Consequently
\begin{equation}
  \frac{A(p)}{A_{\mathrm c}(p)}\longrightarrow\frac12
  \quad(p\downarrow0),
  \qquad
  \frac{A(p)}{A_{\mathrm c}(p)}\longrightarrow1
  \quad(p\uparrow\tfrac12).
  \label{eq:ratiolimits}
\end{equation}
Numerical values increase smoothly between these limits, but no monotonicity
theorem is needed or claimed here.  The rigorous comparison throughout the
interval is \eqref{eq:Ac-comparison-bound}.

\subsection{Closed forms and the Koenigs function}
\label{sec:koenigs}

Several direct searches suggest simple approximations to \(A(p)\).  For example,
a rational symbolic-regression fit obtained from high-precision values was
\begin{equation}
  \widehat A_1(p)=
  \frac{\tfrac12(1-2p)(1+2p)}
  {1+\tfrac12p-\tfrac52p^2-\tfrac65p^3+\tfrac65p^4}.
  \label{eq:A1hat}
\end{equation}
It passes through the two endpoint values, yet
\[
  \frac{\widehat A_1(p)}{2(1-2p)}
  \longrightarrow \frac{10}{11}
  \qquad (p\uparrow\tfrac12),
\]
rather than the limit \(1\) required by \eqref{eq:Aasym}.  Broader tree-based
searches over elementary functions produce many similarly convincing curves
and the same failure of exact constraints.  Near criticality, an error that is
visually negligible on a plot of \(A\) can be substantial after inversion, so
curve fit alone cannot settle the closed-form question.

The structural object behind \(A(p)\) is the linearising coordinate of the
logistic map.  Let \(0<m<1\) and \(f_m(w)=mw(1-w)\).  Koenigs' theorem
\cite{Koenigs1884} gives a unique analytic germ \(\psi_m\) at the attracting
fixed point \(0\) such that
\begin{equation}
  \psi_m(f_m(w))=m\psi_m(w),
  \qquad
  \psi_m(0)=0,
  \qquad
  \psi_m'(0)=1.
  \label{eq:schroder}
\end{equation}
Although this definition is local, the orbit of \(w_0=1/2\) eventually enters
the local neighbourhood.  Equation \eqref{eq:schroder} therefore defines the
required value at \(w_0\) by continuation backwards along that orbit; no claim of
global univalence is required.

\begin{proposition}[Koenigs representation]
For \(0<p<\tfrac12\), with \(m=2p\),
\begin{equation}
  A(p)=2\psi_m\!\left(\frac12\right).
  \label{eq:AeqKoenigs}
\end{equation}
\end{proposition}

\begin{proof}
Let \(w_n=f_m^{\circ n}(1/2)\).  Iterating
\eqref{eq:schroder} gives
\[
  \psi_m\!\left(\frac12\right)
  =\frac{\psi_m(w_n)}{m^n}.
\]
As \(n\to\infty\), \(w_n\to0\) and the normalisation at the origin gives
\(\psi_m(w_n)/w_n\to1\).  Hence
\[
  \psi_m\!\left(\frac12\right)
  =\lim_{n\to\infty}\frac{w_n}{m^n}
  =\frac{A(p)}{2}.
\]
\end{proof}

This representation also locates the limit of elementary methods.  A theorem of
Becker and Bergweiler classifies differentially algebraic solutions of Schr\"oder,
B\"ottcher and Abel equations for rational maps.  In the Schr\"oder case, a
differentially algebraic linearising function can occur only at a repelling fixed
point and then belongs to a restricted family of monomial, Chebyshev or elliptic
types \cite{Becker1995}.  The fixed point here is attracting, because
\(0<m<1\).  Differential algebraicity is preserved under local functional
inversion, so the choice of conjugacy convention does not change the conclusion:
\(w\mapsto\psi_m(w)\) is hypertranscendental, and in particular non-elementary,
for every fixed parameter in the branching range.

The familiar solvable logistic parameters lie outside that range.  At \(m=2\),
the change \(u=1-2w\) converts the iteration into \(u\mapsto u^2\), giving
\begin{equation}
  \psi_2(w)=-\frac12\log(1-2w).
  \label{eq:psi2}
\end{equation}
At \(m=4\), writing \(w=\sin^2\theta\) turns the update into angle doubling and
gives
\begin{equation}
  \psi_4(w)=\bigl(\arcsin\sqrt w\bigr)^2.
  \label{eq:psi4}
\end{equation}
Both concern a repelling origin and serve only as contrasts with \(0<m<1\).

The affine coordinate \(z=m/2-mw\) conjugates the logistic family to
\begin{equation}
  z\longmapsto z^2+c,
  \qquad
  c=\frac{2m-m^2}{4}.
  \label{eq:cofr}
\end{equation}
Hence the branching interval \(m\in(0,1)\) maps to
\(c\in(0,1/4)\), the real centre-to-cusp radius of the main cardioid of the
Mandelbrot set \cite{Milnor2006}.  The parameter map gives useful dynamical
context, but it does not strengthen the closed-form conclusion.

One logical boundary is essential.  Hypertranscendence has been established for
the function \(w\mapsto\psi_m(w)\) with \(m\) fixed.  It does not prove that the
parameter function
\(p\mapsto2\psi_{2p}(1/2)=A(p)\) has no elementary expression.  The supported
conclusion is narrower: no such expression is known; fixed-parameter
linearisation is non-elementary; and numerical searches are consistent with,
but cannot prove, the absence of a simple parameter formula.

\subsection{Certified numerical evaluation}
\label{sec:practical}

The identity \eqref{eq:ratio-is-product} gives a stable calculation and a simple
error certificate.  Let
\[
  R_n=\frac{S_n}{m^n}
  =\prod_{k=0}^{n-1}\left(1-\frac{S_k}{2}\right).
\]
Because \(S_{k+1}<mS_k\),
\(
  \sum_{k=n}^{\infty}S_k/2\leq S_n/[2(1-m)].
\)
The elementary product inequality
\(\prod_k(1-a_k)\geq1-\sum_ka_k\) then gives, whenever the lower endpoint is
positive,
\begin{equation}
  R_n\left(1-\frac{S_n}{2(1-m)}\right)
  \leq A(p)\leq R_n.
  \label{eq:A-certified-interval}
\end{equation}
Thus iteration may stop when the width of this interval meets the required
tolerance.  Accumulating logarithms of the factors avoids underflow.

\begin{table}[htbp]
  \centering
  \caption{Representative values obtained from the monotone ratios and the
  certificate \eqref{eq:A-certified-interval}.  The increasing cost near
  \(p=1/2\) reflects the critical slowdown; the displayed iteration counts use a
  tail target of approximately \(10^{-45}\), not a claim of sixty-digit accuracy.}
  \label{tab:Avals}
  \begin{tabular}{@{}rrrr@{}}
    \toprule
    \(p\) & \(A(p)\) & \(1/A(p)\) & iterations \\
    \midrule
    0.05   & 0.486180336821 & 2.05684994695 & 45 \\
    0.20   & 0.423894537870 & 2.35907734274 & 112 \\
    0.30   & 0.353967177237 & 2.82512070132 & 202 \\
    0.40   & 0.237646658970 & 4.20792787214 & 463 \\
    0.45   & 0.145068746263 & 6.89328353462 & 981 \\
    0.49   & 0.036394897530 & 27.4763790497 & 5\,125 \\
    0.499  & 0.003944476201 & 253.519085680 & 51\,750 \\
    0.4999 & 0.000399257804 & 2504.64734704 & 518\,021 \\
    \bottomrule
  \end{tabular}
\end{table}

For exploratory work extremely close to criticality,
\(2(1-2p)\) is an inexpensive leading approximation.  It should not be joined to
the product at an arbitrary fixed crossover when precision matters: the empirical
logarithmic correction makes its relative error decay slowly.  A reliable hybrid
scheme uses \eqref{eq:A-certified-interval} wherever iteration is feasible and
switches to the asymptotic only after its error has been calibrated at the required
parameter and tolerance.

\FloatBarrier
